% §5.8 The general linear row along the chain --- the statement and the proof of thm:wedge-axis
% (Q-0290). Q-0309 (fix group 2, review flag 10): the formal statement moved here from 5.3
% (05-table.tex), where it stood before its vocabulary (C_alpha and its order, 5.6; the order of a
% cone on PSD(d), 5.4; the stable laws and the dilation ratio, 5.6); 5.3 keeps the table and an
% informal statement. Included from 05-landscape.tex; the conventions are in that file's header.

\subsection{The general linear row along the chain}\label{sec:wedge-axis-proof}

% Writing-rules pass, group s47 (rule 10): the opening states the question in the reader's terms
% and the answer by clause; the trust-base sentence is gone (section 1.1 owns it); the paragraph on
% the two reasons for the stop at alpha = 2 moved before the theorem.
A practitioner who keeps covariance under the whole general linear group and wants a cascade across
shapes with a kernel other than the Gaussian has to give up part of the Loewner order, and the part kept decides how fast the shape may change from one
stage to the next. Along the Loewner order the anisotropy gained in a step from volume scale $t_1$
to $t_2$ is at most $(t_2/t_1)^d$ (Lemma~\ref{lem:loewner-scale-shape}(3)); along the stable order
of index $\alpha$ it is at most $(t_2/t_1)^{d\alpha/(d+\alpha)}$
(Proposition~\ref{prop:stable-step-cone}(3)); along the rays the shape does not change. So the
question is how large an order a non-Gaussian kernel can keep, and which kernels keep it.
Theorem~\ref{thm:wedge-axis} answers it for the orders generated by the invariant cones of
Theorem~\ref{thm:cone-orders}(4). A cone not contained in $\PSD(d)$ carries no family, and a cone in
$\PSD(d)$ that contains a nonzero singular matrix, the direction of a rank-deficient step, leaves
the affine Gaussian scale spaces alone. For a cone whose nonzero members are positive definite, one
number $\beta(C)$ decides: a family whose kernel at $I$ is not Gaussian exists exactly when
$\beta(C) < 2$, and for $0 < \gamma < 2$ the orbit of the isotropic $\gamma$-stable law is a family
exactly when $\gamma \ge \beta(C)$ (clause (2)). The stable kernels are the familiar non-Gaussian
ones, those that the $\alpha$-scale spaces of Pauwels et al.\ and of Duits et al.\ admit where their
kernels are positive
(\sscite{pauwels1995extended}, p.~694; \sscite{duits2004axioms}, p.~277). Along the chain $C_\alpha$
the orbit of each keeps a cascade up to the order of its own index, and from $\alpha = 2$ on no
non-Gaussian kernel keeps one (clause (3)).

The theorem assembles the general linear cells of
\S\S\ref{sec:invariant-wedges}--\ref{sec:affine-stable-cone} along the chain of cones $C_\alpha$. For an
$\Orth(d)$-invariant closed convex cone $C \subseteq \Sym(d)$
with $I \in C$ and $C \cap (-C) = \{0\}$, the order $\preceq^+_C$ on $\PSD(d)$ relates two definite
matrices when a curve whose relative velocity lies in $C$ at almost every time joins them, and adds the
origin and the rays (Theorem~\ref{thm:cone-orders}(4)). The cones
$C_\alpha := \{X \in \PSD(d) : (d+\alpha)\,\lambda_{\min}(X) \ge \tr X\}$, $0 \le \alpha < \infty$, and
$C_\infty := \PSD(d)$ form a chain from the rays ($\alpha = 0$) through the stable orders
($0 < \alpha < 2$) and the cones $C_\alpha$, $2 \le \alpha < \infty$, to the Loewner order
($\alpha = \infty$), and $\preceq_\alpha$ is the order of
$C_\alpha$ on $\PD(d)$ (Proposition~\ref{prop:stable-order-generated}). For $0 < \alpha < 2$ it is the stable
order, and on $\PSD(d)$ the order $\preceq^+_{C_\alpha}$ is the $\le_\alpha$ of
Proposition~\ref{prop:similarity-stable-order-families}(1) (Table~\ref{tab:notation}). Here $S_\gamma(I)$ is the rotation-invariant
$\gamma$-stable law, with transform $e^{-|\omega|^\gamma}$ (Lemma~\ref{lem:elliptical-stable}).

Non-Gaussian families stop at $\alpha = 2$, below the Loewner order, for two reasons, one for the cones
that reach the boundary of $\PSD(d)$ and one for those that do not. When a cone reaches the boundary,
whose nonzero points are the nonzero singular matrices, a single rank-deficient step forces the Gaussian
(Theorem~\ref{thm:rank-deficient-step}), and no $C_\alpha$ with $\alpha < \infty$ reaches that boundary.
Between $C_2$ and the boundary the anisotropy budget of Proposition~\ref{prop:anisotropy-budget}
excludes every non-Gaussian law, at $\alpha = 2$ itself included, since the stable cones fill only the
open part of $C_2$ (Theorem~\ref{thm:wedge-axis}(1)).

% shared with blueprint thm:wedge-axis+e9a765e87583
\begin{theorem}[the general linear row along the chain of cones]\label{thm:wedge-axis}
Let $d \ge 2$, and let $C_\alpha$ and $\preceq_{C_\alpha}$ be as in Proposition~\ref{prop:stable-order-generated} and
$\preceq^+_C$ as in Theorem~\ref{thm:cone-orders}(4).
\begin{enumerate}
\item \emph{(The chain.)} For $0 \le \alpha < \beta \le \infty$, $C_\alpha \subsetneq C_\beta$ and
$\preceq_{C_\alpha} \subsetneq \preceq_{C_\beta}$. For $\alpha < \infty$ every nonzero member of $C_\alpha$ is positive
definite, while $C_\infty = \PSD(d)$ contains every nonzero singular positive semidefinite matrix. The union of
the $C_\alpha$, $0 \le \alpha < 2$, is $\{0\} \cup \{X \in \Sym(d) : (d+2)\lambda_{\min}(X) > \tr X\}$, a proper
subset of $C_2$.
\item \emph{(Which cones carry general linear families.)} Let $C \subseteq \Sym(d)$ be an $\Orth(d)$-invariant
closed convex cone with $I \in C$ and $C \cap (-C) = \{0\}$, and consider the
$(\GL(d), \preceq^+_C)$-cascade families of kernels (Definition~\ref{def:gw-cascade-family}(2)).
\begin{enumerate}
\item[(a)] If $C \not\subseteq \PSD(d)$, there is none.
\item[(b)] If $C \subseteq \PSD(d)$ and $C$ contains a nonzero singular matrix, they are exactly the
restrictions to the pairs $\Sigma \preceq^+_C \Sigma'$ of the affine Gaussian scale spaces
\ref{eq:affine-gaussian}, $b > 0$.
\item[(c)] Otherwise every nonzero member of $C$ is positive definite, the largest value $\beta(C)$ of
$\tr X/\lambda_{\min}(X) - d$ over $X \in C \setminus \{0\}$ is attained, and $\beta(C) \in [0, \infty)$ is the
least $\alpha$ with $C \subseteq C_\alpha$. For $0 < \gamma < 2$, $\Sigma \mapsto (\Sigma^{1/2})_*S_\gamma(I)$
is the family of kernels from scale zero of a $(\GL(d), \preceq^+_C)$-cascade family if and only if
$\gamma \ge \beta(C)$. A $(\GL(d), \preceq^+_C)$-cascade family whose $\mu_I$ is not Gaussian exists if and only
if $\beta(C) < 2$; if $\beta(C) \ge 2$, the families are, as in (b), exactly the restrictions to the pairs
$\Sigma \preceq^+_C \Sigma'$ of the affine Gaussian scale spaces \ref{eq:affine-gaussian}, $b > 0$.
\end{enumerate}
\item \emph{(The row along the chain.)} For $0 \le \alpha \le \infty$ let $\mathfrak F_\alpha$ be the set of
laws $\mu_I$ of $(\GL(d), \preceq^+_{C_\alpha})$-cascade families of kernels.
\begin{enumerate}
\item[(a)] $\mathfrak F_0$ is the set of rotation-invariant self-decomposable laws other than $\delta_0$.
\item[(b)] $\mathfrak F_\beta \subseteq \mathfrak F_\alpha$ for $\alpha \le \beta$, strictly when
$\alpha < \beta \le 2$.
\item[(c)] For $0 < \alpha < 2$, $\mathfrak F_\alpha$ contains $S_\gamma(I)$, $0 < \gamma < 2$, if and only
if $\gamma \ge \alpha$, and contains the law $\mu_I$ of Proposition~\ref{prop:orbit-step-cone}(1) for every profile whose
dilation ratio has $R(\lambda) \ge \lambda^\alpha$ for $\lambda \ge 1$.
\item[(d)] For $2 \le \alpha \le \infty$, $\mathfrak F_\alpha$ is the set of centred Gaussian laws with
covariance $2bI$, $b > 0$.
\end{enumerate}
\end{enumerate}
\end{theorem}

So the general linear row shrinks from the orbits of all rotation-invariant self-decomposable laws other
than $\delta_0$ to the Gaussian, and for the stable
orbits the stable index decides exactly where on the chain they stop. For the orbit of
every profile, with any Gaussian part, the converse of the profile clause of (3)(c) holds as well
(Proposition~\ref{prop:wedge-axis-profiles} below).

% Round-1 review (Gemini), the author's decision D3 (2026-11-03): the body sketch of the proof is
% replaced by its idea and a pointer; Appendix C prints the complete proof of record, matched by its
% heading.
The proof reads clauses (2) and (3) off the steps $I \preceq^+_C e^X$, $X \in C$, and the increment that a
family gives each of them. A negative eigenvalue of $X$ leaves no family, and a singular $X \ne 0$ makes
the step rank-deficient, which forces the Gaussian (Theorem~\ref{thm:rank-deficient-step}). For a cone of
definite matrices the stable cone of Proposition~\ref{prop:stable-step-cone}(1) and the anisotropy budget
of Proposition~\ref{prop:anisotropy-budget}(1) give the threshold $\beta(C)$, and clause (3) is clause (2)
for the cones $C_\alpha$, which have $\beta(C_\alpha) = \alpha$. The complete proof, with complete proofs of
the statements it spends, is in Appendix~\ref{app:wedge-axis}.

% Round 3 (report E, section 3.1; the author's decision D10, PLAN-review-response-paper.md R3.0): the
% converse of the profile clause, the cone orders, and the whole cell with C5.
The increments of a family along these orders are infinitely divisible (Theorem~\ref{thm:cascade-infinitely-divisible}(2)),
and an infinitely divisible increment orders the \Levy{} densities
(Proposition~\ref{prop:similarity-orbit-family}(1)). So for families the ordered-density condition of
Proposition~\ref{prop:orbit-step-cone} is necessary, and the stable rung of a profile, its lower exponent, decides
where its orbit stops on the chain. With C5 this describes the cells for every finite $\alpha > 0$ completely.
From $\alpha = 2$ on no profile is left, which gives the laws of Theorem~\ref{thm:wedge-axis}(3)(d) for finite
$\alpha$ without C3.

% CHANGED the author's decision D14 (2026-11-04; PLAN-review-response-paper.md R3.0): clause (1) holds
% for every profile; the statement's k is no longer right-continuous, and the proof gains "(1), the
% right-continuous modification". Clause (2) keeps the word. Re-transcribed from the blueprint (adr/0004).
% CHANGED the author's decision D16 (2026-11-04; same plan): clause (2) drops alpha < 2 and beta(C) < 2,
% which the proofs did not read, and says what is left from alpha = 2 on. Re-transcribed (adr/0004).
% shared with blueprint prop:wedge-axis-profiles
\begin{proposition}[the general linear row along the chain: the profiles, and the whole cell]\label{prop:wedge-axis-profiles}
Let $d \ge 2$, and let $C_\alpha$, $\preceq^+_C$, the number $\beta(C)$ and $\mathfrak F_\alpha$ be as in
Theorem~\ref{thm:wedge-axis}. For a cone $C$ as in Theorem~\ref{thm:wedge-axis}(2) with $C \subseteq \PSD(d)$ and no nonzero
singular member, the case (2)(c) there, write $\mathfrak F(C)$ for the set of laws $\mu_I$ of
$(\GL(d), \preceq^+_C)$-cascade families of kernels, so that $\mathfrak F_\alpha = \mathfrak F(C_\alpha)$ for
$0 \le \alpha < \infty$. Let $k$ be a profile as in Proposition~\ref{prop:orbit-step-cone}, with dilation
ratio $R$ and lower exponent $\beta_-(k)$ (Lemma~\ref{lem:dilation-ratio}(3)), let $\nu_I$ be its law
(Proposition~\ref{prop:orbit-step-cone}(1)), and for $a \ge 0$ write $N(0, aI)$ for the centred Gaussian law with
covariance $aI$.
\begin{enumerate}
\item \emph{(The profiles.)} Let $C$ be such a cone, with $\beta(C) > 0$. Then $N(0, aI) * \nu_I \in \mathfrak F(C)$
if and only if $\beta(C) \le \beta_-(k)$. In particular, for $0 < \alpha < 2$, $N(0, aI) * \nu_I \in \mathfrak F_\alpha$
if and only if $R(\lambda) \ge \lambda^\alpha$ for every $\lambda \ge 1$, and then
$\Sigma \mapsto N(0, a\Sigma) * (\Sigma^{1/2})_*\nu_I$, $\Sigma \in \PSD(d)$, is the family of kernels from scale
zero of a $(\GL(d), \preceq^+_{C_\alpha})$-cascade family of kernels.
\item \emph{(The whole cell.)} For $0 < \alpha < \infty$, $\mathfrak F_\alpha$ consists of the laws $N(0, aI)$, $a > 0$,
and the laws $N(0, aI) * \nu_I$, $a \ge 0$, of the right-continuous profiles $k$ with $\beta_-(k) \ge \alpha$.
From $\alpha = 2$ on there is no such profile, so $\mathfrak F_\alpha$ consists of the laws $N(0, aI)$, $a > 0$.
For a cone $C$ as above, $\mathfrak F(C) = \mathfrak F_{\beta(C)}$: the laws $\mu_I$ of the families
along $\preceq^+_C$ depend on $C$ only through $\beta(C)$.
\end{enumerate}
\end{proposition}

\begin{proof}
\emph{The orders.} For $0 < \gamma < 2$, $\preceq^+_{C_\gamma}$ is the order $\le_\gamma$ of
Proposition~\ref{prop:similarity-stable-order-families}(1): on $\PD(d)$ both are $\preceq_\gamma$
(Proposition~\ref{prop:stable-order-generated}(1)), and both add the origin and the rays. Let $C$ be a cone as in the
statement. By Theorem~\ref{thm:wedge-axis}(2)(c) the largest value $\beta(C)$ of $\tr X/\lambda_{\min}(X) - d$ over
$X \in C \setminus \{0\}$ is attained and $C \subseteq C_{\beta(C)}$. So $\preceq^+_C$ is contained in
$\preceq^+_{C_{\beta(C)}}$, a $C$-causal curve being $C_{\beta(C)}$-causal, and a
$(\GL(d), \preceq^+_{C_{\beta(C)}})$-cascade family restricts to a $(\GL(d), \preceq^+_C)$-cascade family with the
same $\mu_I$ (the proof of Theorem~\ref{thm:wedge-axis}, the orders). For $X \in C$ and $t \ge 0$ the curve
$\zeta(s) := e^{stX}$, $s \in [0,1]$, is a chain from $I$ to $e^{tX}$ for $\preceq^+_C$ as in
Theorem~\ref{thm:cascade-infinitely-divisible}(2). It is continuous, and for $s \le s'$, $I \preceq^+_C e^{(s'-s)tX}$
as $(s'-s)tX \in C$ (the proof of Theorem~\ref{thm:wedge-axis}, the orders). The matrix $A := e^{stX/2}$ carries
$I$ to $\zeta(s)$ and $e^{(s'-s)tX}$ to $\zeta(s')$, so $\zeta(s) \preceq^+_C \zeta(s')$ by the
$\GL(d)$-invariance of the order.

\emph{(1), sufficiency.} Let $\beta(C) \le \beta_-(k)$ and put $\gamma := \beta(C)$. Then $0 < \gamma < 2$, as
$\beta_-(k) < 2$, and $R(\lambda) \ge \lambda^\gamma$ for every $\lambda \ge 1$ (Lemma~\ref{lem:dilation-ratio}(3)). By
Proposition~\ref{prop:similarity-stable-order-families}(2) the laws $\nu_\Sigma := (\Sigma^{1/2})_*\nu_I$, $\Sigma \in \PSD(d)$,
are the kernels from scale zero of a $(\GL(d), \le_\gamma)$-cascade family $(\nu_{\Sigma,\Sigma'})$. Put
$\mu_\Sigma := N(0, a\Sigma) * \nu_\Sigma$. Its transform $e^{-a\langle\omega,\Sigma\omega\rangle/2}\,\widehat\nu_\Sigma(\omega)$
has no zero, as $\widehat\nu_\Sigma(\omega) = e^{-\psi_I(\Sigma^{1/2}\omega)}$ (Proposition~\ref{prop:orbit-step-cone}(1)). The
conditions (a)--(d) of Lemma~\ref{lem:gw-kernels-from-origin}(2) hold for $G = \GL(d)$ and the order $\le_\gamma$, as
they hold for each factor.
\begin{itemize}
\item (a): $\mu_0 = \delta_0$, and $\mu_I \ne \delta_0$, as $\widehat\mu_I = e^{-a|\omega|^2/2 - \psi_I}$ and $\psi_I$
is not identically $0$, $k$ being positive near $0$.
\item (b): $A_*N(0, a\Sigma) = N(0, aA\Sigma A^{\tp})$ and $A_*\nu_\Sigma = \nu_{A\cdot\Sigma}$ for $A \in \GL(d)$.
\item (c): both transforms are continuous in $\Sigma$.
\item (d): for $\Sigma \le_\gamma \Sigma'$, $\Sigma' - \Sigma \ge 0$ (Proposition~\ref{prop:similarity-stable-order-families}(1)), so
$\rho := N(0, a(\Sigma' - \Sigma)) * \nu_{\Sigma,\Sigma'}$ is a probability measure, and $\mu_{\Sigma'} = \mu_\Sigma * \rho$.
\end{itemize}
By Lemma~\ref{lem:gw-kernels-from-origin}(3) the $\mu_\Sigma$ are the kernels from scale zero of a
$(\GL(d), \le_\gamma)$-cascade family, whose law at $I$ is $N(0, aI) * \nu_I$, and it restricts to $\preceq^+_C$.

\emph{(1), the right-continuous modification.} Put $\tilde k(r) := k(r{+}) = \lim_{s \downarrow r}k(s)$ for $r > 0$.
It is nonincreasing and right-continuous, and $\tilde k \le k$, so $\tilde k$ satisfies the
integrability conditions that $k$ does. If $k(r_0) > 0$, then $\tilde k \ge k(r_0)$ on
$(0, r_0)$, so $\tilde k$ is a profile as in Proposition~\ref{prop:orbit-step-cone}. A monotone function has at most
countably many discontinuities, so $\tilde k = k$ off a countable subset of $(0, \infty)$, and the densities
$k(|x|)|x|^{-d}$ and $\tilde k(|x|)|x|^{-d}$ agree off the set of the $x$ with $|x|$ in that countable set,
a null set of $\RR^d$. So the two profiles have the same
$\psi_I$ and the same law $\nu_I$. They also have the same lower exponent. Let $\tilde R$ be the dilation ratio of
$\tilde k$ and $\lambda \ge 1$. Every $s > 0$ has $k(s/\lambda) \ge R(\lambda)k(s)$, by the definition of $R$ when
$k(s) > 0$ and trivially otherwise. Letting $s$ decrease to $r$ gives $\tilde k(r/\lambda) \ge R(\lambda)\tilde k(r)$,
so $\tilde R(\lambda) \ge R(\lambda)$. Conversely, let $\alpha \ge 0$ with $\tilde R(\tau) \ge \tau^\alpha$ for every
$\tau \ge 1$, let $\lambda > 1$, and let $k(r) > 0$. For $r/\lambda < s < r$, $\tilde k(s) \ge k(r) > 0$, as
$k \ge k(r)$ on $(s, r)$, and $\tau := s\lambda/r > 1$. So
$k(r/\lambda) \ge \tilde k(r/\lambda) = \tilde k(s/\tau) \ge \tau^\alpha\tilde k(s) \ge (s\lambda/r)^\alpha k(r)$, and
letting $s$ increase to $r$ gives $k(r/\lambda) \ge \lambda^\alpha k(r)$. So $R(\lambda) \ge \lambda^\alpha$, which
holds at $\lambda = 1$ as well. By Lemma~\ref{lem:dilation-ratio}(3), applied to $k$ and to $\tilde k$, the numbers
$\alpha \ge 0$ with $\alpha \le \beta_-(k)$ are those with $\alpha \le \beta_-(\tilde k)$, and
$\beta_-(k) = \beta_-(\tilde k)$. Both sides of the equivalence in (1) are therefore the same for $k$ and $\tilde k$.
In the necessity $k$ may be replaced by $\tilde k$, and so taken right-continuous.

\emph{(1), necessity.} By the modification we may take $k$ right-continuous. Let $\mu_I := N(0, aI) * \nu_I$ be the law at $I$ of a $(\GL(d), \preceq^+_C)$-cascade
family $(\mu_{\Sigma,\Sigma'})$. For $\Sigma \in \PD(d)$, (K3) with $A = \Sigma^{1/2}$ at the pair $(0, I)$ gives
$\mu_\Sigma = (\Sigma^{1/2})_*\mu_I = N(0, a\Sigma) * (\Sigma^{1/2})_*\nu_I$. Let $X \in C \setminus \{0\}$ attain
$\beta(C)$, so that $\tr X = (d + \beta(C))\lambda$ with $\lambda := \lambda_{\min}(X) > 0$, $X$ being positive definite. Let $t > 0$ and $M := e^{tX}$. By the chain of the first paragraph and
Theorem~\ref{thm:cascade-infinitely-divisible}(2) the increment $\rho := \mu_{I,M}$ is infinitely divisible, and (K1) at
$0 \preceq^+_C I \preceq^+_C M$ gives $\mu_M = \mu_I * \rho$. So Proposition~\ref{prop:similarity-orbit-family}(1), with $h$ the
identity and the Gaussian covariance $\Sigma \mapsto a\Sigma$, puts the step $I \to M$ in the step set of
Proposition~\ref{prop:orbit-step-cone}(2): $\det M \le \lambda_{\min}(M)^dR(\lambda_{\min}(M)^{1/2})^2$. Here
$\lambda_{\min}(M) = e^{t\lambda}$ and $\det M = e^{t\tr X} = e^{t(d+\beta(C))\lambda}$, so with $u := e^{t\lambda/2}$ the
inequality reads $u^{2(d+\beta(C))} \le u^{2d}R(u)^2$, that is $R(u) \ge u^{\beta(C)}$. As $t$ ranges over
$(0, \infty)$, $u$ ranges over $(1, \infty)$, and $R(1) = 1$. So $\beta(C) \le \beta_-(k)$ by
Lemma~\ref{lem:dilation-ratio}(3). For $C = C_\alpha$ one may take $X = \diag(1, \dots, 1, 1 + \alpha)$, and $M$ is the
step of the proof of Lemma~\ref{lem:dilation-ratio}(4).

\emph{(1), the stable orders.} For $0 < \alpha < 2$, $C_\alpha$ is a cone as in the statement with
$\beta(C_\alpha) = \alpha$ (the proof of Theorem~\ref{thm:wedge-axis}, (3)), and $\mathfrak F_\alpha = \mathfrak F(C_\alpha)$. By
Lemma~\ref{lem:dilation-ratio}(3), $\alpha \le \beta_-(k)$ if and only if $R(\lambda) \ge \lambda^\alpha$ for every
$\lambda \ge 1$, and the family is the one built in the sufficiency with $\gamma = \alpha$, along
$\le_\alpha = \preceq^+_{C_\alpha}$.

\emph{(2), the form of $\mu_I$.} Let $C$ be a cone as in the statement and $\mu_I \in \mathfrak F(C)$. The ray order
is contained in $\preceq^+_C$, so the family restricts to a $(\GL(d), \text{rays})$-cascade family with the same
$\mu_I$, and by Proposition~\ref{prop:gw-ray-families}(3) and (1), $\mu_I$ is rotation invariant, self-decomposable and not
$\delta_0$, and $\widehat\mu_I$ has no zero. The pair $0 \preceq^+_C I$ is joined by the chain $s \mapsto sI$, so
$\mu_I = \mu_{0,I}$ is infinitely divisible (Theorem~\ref{thm:cascade-infinitely-divisible}(2)). Its transform is radial,
and real as $\mu_I$ is invariant under $-I$. Having no zero, it is $\widehat\mu_I(\omega) = e^{-F(|\omega|)}$ with
$F(r) := -\log\widehat\mu_I(re_1) \ge 0$. By \ledger{A3}(i), $\mu_I$ has a generating triplet $(A, \nu, \gamma)$.
The real part of $e^{i\theta} - 1 - i\theta\mathbf 1_{\{|x| \le 1\}}$ is $\cos\theta - 1$, so taking moduli in the
representation,
\[
  F(|\omega|) = \tfrac12\langle\omega, A\omega\rangle + \int_{\RR^d}\bigl(1 - \cos\langle\omega,x\rangle\bigr)\,\nu(dx) ,
  \qquad \omega \in \RR^d .
\]
As $\mu_I$ is self-decomposable, the polar criterion \ledger{A2} gives
$\nu(B) = \int_S\lambda(d\xi)\int_0^\infty\mathbf 1_B(r\xi)\,k_\xi(r)\,dr/r$ for Borel $B$, with $\lambda$ a finite
measure on $S = S^{d-1}$ and $k_\xi(r) \ge 0$ measurable in $\xi$ and nonincreasing in $r$. A nonincreasing $k_\xi$
differs from its right limit $k_\xi(r+) = \sup\{k_\xi(q) : q \in \mathbb Q,\ q > r\}$ at countably many $r$ only,
so replacing $k_\xi(r)$ by $k_\xi(r+)$ changes no inner integral, and it makes the kernel jointly measurable in
$(\xi, r)$, as a countable supremum of measurable functions. Let $\sigma_{d-1}$ be the surface measure of $S$, with total
$|S^{d-1}| = \sigma_{d-1}(S)$, and put
\[
  k(r) := \frac{1}{|S^{d-1}|}\int_S k_\xi(r+)\,\lambda(d\xi) \in [0, \infty] , \qquad r > 0 .
\]
It is nonincreasing. Exchanging the integrals (the integrands are nonnegative) and passing to polar coordinates
$x = r\eta$, $dx = r^{d-1}\,dr\,\sigma_{d-1}(d\eta)$,
\[
  \int_{\RR^d}(|x|^2 \wedge 1)\,\nu(dx) = |S^{d-1}|\int_0^\infty(r^2 \wedge 1)\,k(r)\,\frac{dr}{r}
  = \int_{\RR^d}(|x|^2 \wedge 1)\,k(|x|)\,|x|^{-d}\,dx ,
\]
which is finite by \ledger{A3}(i). So $k$ is finite: a nonincreasing function infinite at some $r_0$ is infinite
on $(0, r_0]$, and the integral would diverge. The left side of the display for $F$ depends on $\omega$ only
through $|\omega|$, so it equals the mean of the right side at the points $|\omega|\theta$, $\theta \in S$, against
$\sigma_{d-1}/|S^{d-1}|$. The Gaussian term has the mean $\tfrac12 a|\omega|^2$ with $a := \tr A/d \ge 0$, as
$\int_S\theta_i\theta_j\,\sigma_{d-1}(d\theta) = \delta_{ij}|S^{d-1}|/d$. For the other term let
$c(s) := |S^{d-1}|^{-1}\int_S(1 - \cos(s\langle\theta,\xi\rangle))\,\sigma_{d-1}(d\theta)$ for $s \ge 0$ and a unit vector $\xi$. It does not depend on $\xi$, since $\sigma_{d-1}$ is invariant under a rotation carrying one unit vector to
another. Exchanging the integrals as before and then passing to polar coordinates, with
$c(r|\omega|) = |S^{d-1}|^{-1}\int_S(1 - \cos\langle\omega, r\eta\rangle)\,\sigma_{d-1}(d\eta)$ (take $\xi = \omega/|\omega|$
when $\omega \ne 0$),
\[
  F(|\omega|) = \tfrac12 a|\omega|^2 + |S^{d-1}|\int_0^\infty c(r|\omega|)\,k(r)\,\frac{dr}{r}
  = \tfrac12 a|\omega|^2 + \int_{\RR^d}\bigl(1 - \cos\langle\omega,x\rangle\bigr)\,k(|x|)\,|x|^{-d}\,dx ,
  \qquad \omega \in \RR^d .
\]
If $k \equiv 0$, then $\widehat\mu_I(\omega) = e^{-a|\omega|^2/2}$ and
$\mu_I = N(0, aI)$, with $a > 0$ as $\mu_I \ne \delta_0$. Otherwise $k$ is a profile as in
Proposition~\ref{prop:orbit-step-cone}, and so is its right-continuous modification $\tilde k$, with the same exponent (the paragraph
on the modification). So $\widehat\mu_I = e^{-a|\omega|^2/2}\,\widehat\nu_I$ for the law $\nu_I$ of $\tilde k$, and
$\mu_I = N(0, aI) * \nu_I$ by the uniqueness of the transform.

\emph{(2), the equality.} If $\beta(C) = 0$, then $C \subseteq C_0 = \{cI : c \ge 0\}$
(Proposition~\ref{prop:stable-order-generated}(1)), and $C_0 \subseteq C$ as $I \in C$. So $C = C_0$, and there is nothing to prove. Let $\beta(C) > 0$. By the first paragraph, $\mathfrak F_{\beta(C)} \subseteq \mathfrak F(C)$.
Conversely let $\mu_I \in \mathfrak F(C)$, of one of the two forms above. A law $N(0, aI)$, $a > 0$, is the law
at $I$ of the affine Gaussian scale space with constant $a/2$, an affine cascade family
(Theorem~\ref{thm:affine-rigidity}), and so of its restriction to the pairs of $\preceq^+_{C_{\beta(C)}}$, which are
Loewner pairs, $C_{\beta(C)}$ being contained in $\PSD(d)$ (the proof of Theorem~\ref{thm:wedge-axis}, the orders). So it lies in $\mathfrak F_{\beta(C)}$. A law $N(0, aI) * \nu_I$ has $\beta(C) \le \beta_-(k)$ by (1) for $C$, and so
lies in $\mathfrak F(C_{\beta(C)}) = \mathfrak F_{\beta(C)}$ by (1) for $C_{\beta(C)}$, whose number is $\beta(C)$. Hence
$\mathfrak F(C) = \mathfrak F_{\beta(C)}$. For $C = C_\alpha$, $0 < \alpha < \infty$, a cone as in the statement with
$\beta(C_\alpha) = \alpha$ (the proof of Theorem~\ref{thm:wedge-axis}, (3)), the same two cases describe
$\mathfrak F_\alpha$: each member has one of the two forms, each $N(0, aI)$ with $a > 0$ lies in it, and
$N(0, aI) * \nu_I$ lies in it exactly when $\beta_-(k) \ge \alpha$, by (1). For $\alpha \ge 2$ the second case does
not occur, as $\beta_-(k) < 2$ for every profile (Lemma~\ref{lem:dilation-ratio}(3)).
\end{proof}

The single-step converse is not used, and it stays open. A positive increment that is not known to be
infinitely divisible may fail to order the densities, as the computation of
Remark~\ref{rem:gaussian-matern-step} suggests.

